
\documentclass[aps,preprint]{revtex4}
%%%%%%%%%%%%%%%%%%%%%%%%%%%%%%%%%%%%%%%%%%%%%%%%%%%%%%%%%%%%%%%%%%%%%%%%%%%%%%%%%%%%%%%%%%%%%%%%%%%%%%%%%%%%%%%%%%%%%%%%%%%%%%%%%%%%%%%%%%%%%%%%%%%%%%%%%%%%%%%%%%%%%%%%%%%%%%%%%%%%%%%%%%%%%%%%%%%%%%%%%%%%%%%%%%%%%%%%%%%%%%%%%%%%%%%%%%%%%%%%%%%%%%%%%%%%
\usepackage{amsfonts}

%TCIDATA{OutputFilter=LATEX.DLL}
%TCIDATA{Version=5.50.0.2960}
%TCIDATA{<META NAME="SaveForMode" CONTENT="1">}
%TCIDATA{BibliographyScheme=Manual}
%TCIDATA{Created=Sunday, December 30, 2001 12:27:31}
%TCIDATA{LastRevised=Saturday, July 03, 2010 09:57:04}
%TCIDATA{<META NAME="GraphicsSave" CONTENT="32">}
%TCIDATA{<META NAME="DocumentShell" CONTENT="Articles\SW\REVTeX 4">}
%TCIDATA{Language=American English}
%TCIDATA{CSTFile=revtex4.cst}

\newtheorem{theorem}{Theorem}
\newtheorem{acknowledgement}[theorem]{Acknowledgement}
\newtheorem{algorithm}[theorem]{Algorithm}
\newtheorem{axiom}[theorem]{Axiom}
\newtheorem{claim}[theorem]{Claim}
\newtheorem{conclusion}[theorem]{Conclusion}
\newtheorem{condition}[theorem]{Condition}
\newtheorem{conjecture}[theorem]{Conjecture}
\newtheorem{corollary}[theorem]{Corollary}
\newtheorem{criterion}[theorem]{Criterion}
\newtheorem{definition}[theorem]{Definition}
\newtheorem{example}[theorem]{Example}
\newtheorem{exercise}[theorem]{Exercise}
\newtheorem{lemma}[theorem]{Lemma}
\newtheorem{notation}[theorem]{Notation}
\newtheorem{problem}[theorem]{Problem}
\newtheorem{proposition}[theorem]{Proposition}
\newtheorem{remark}[theorem]{Remark}
\newtheorem{solution}[theorem]{Solution}
\newtheorem{summary}[theorem]{Summary}
\newenvironment{proof}[1][Proof]{\noindent\textbf{#1.} }{\ \rule{0.5em}{0.5em}}
\input{tcilatex}
\begin{document}

\preprint{Zuminos' Theorem}
\title[Zuminos' Theorem]{Zuminos' Theorem}
\author{Brian Weiner}
\affiliation{PSU}
\date{06/17/2003}
\startpage{1}
\endpage{2}
\maketitle
\tableofcontents

\section{Canonical Coordinates}

\begin{theorem}
If $\mathbf{c}$ is a $2s\times 2s$ complex antisymmetric matrix then there
exists a unitary matrix $\mathbb{T}$ such that 
\begin{equation}
\mathbb{T}\mathbf{c}\mathbb{T}^{t}\mathbb{=}\left( 
\begin{array}{cc}
\mathbf{0} & \mathbf{\zeta } \\ 
-\mathbf{\zeta } & \mathbf{0}%
\end{array}%
\right)
\end{equation}%
where 
\begin{equation}
\mathbf{\zeta =}\left( 
\begin{array}{ccc}
\zeta _{1} & 0 & 0 \\ 
0 & \ddots & 0 \\ 
0 & 0 & \zeta _{s}%
\end{array}%
\right) ;\qquad \zeta _{j}\in \mathbb{C},\quad 1\leq j\leq s
\end{equation}

\begin{proof}
We first note some relationships implied by the antisymmetry of $\mathbf{c}$%
\begin{equation}
\mathbf{c}^{\dagger }\mathbf{c\,=-c}^{\ast }\mathbf{c}=\left( \mathbf{cc}%
^{\dagger }\right) ^{t}=\mathbf{\,}\left( \mathbf{cc}^{\dagger }\right)
^{\ast }  \label{cc}
\end{equation}%
(so $\mathbf{c}$ is not a normal operator in general) and one by the
associativity of matrix product 
\begin{equation}
\mathbf{c}\left( \mathbf{c}^{\dagger }\mathbf{c}\right) =\left( \mathbf{cc}%
^{\dagger }\right) \mathbf{c}
\end{equation}%
which by using eq. $\left( \ref{cc}\right) $ gives 
\begin{equation}
\left( \mathbf{cc}^{\dagger }\right) \mathbf{c=c}\left( \mathbf{cc}^{\dagger
}\right) ^{\ast }  \label{cr}
\end{equation}%
Then we introduce $\mathbf{c}_{U}\mathbf{c}_{U}^{\dagger }$ through the
definition 
\begin{eqnarray}
\mathbf{c} &=&\mathbb{U}\mathbf{c}_{U}\mathbb{U}^{t} \\
\mathbf{c}_{U} &=&\mathbb{U}^{\dagger }\mathbf{c}\mathbb{U}^{\ast }
\end{eqnarray}%
where $\mathbb{U}$ is any unitary matrix and obtain that 
\begin{eqnarray}
\mathbf{cc}^{\dagger } &=&\mathbb{U}\mathbf{c}_{U}\mathbb{U}^{t}\mathbb{U}%
^{\ast }\mathbf{c}_{U}^{\dagger }\mathbb{U}^{\dagger }=\mathbb{U}\mathbf{c}%
_{U}\mathbf{c}_{U}^{\dagger }\mathbb{U}^{\dagger }  \nonumber \\
\mathbf{c}_{U}\mathbf{c}_{U}^{\dagger } &=&\mathbb{U}^{\dagger }\mathbf{cc}%
^{\dagger }\mathbb{U}
\end{eqnarray}%
and as $\mathbf{c}_{U}$ is also antisymmetric 
\begin{equation}
\mathbf{c}_{U}\mathbf{c}_{U}^{\dagger }=-\mathbf{c}_{U}\mathbf{c}_{U}^{\ast
}=\left( \mathbf{c}_{U}^{\dagger }\mathbf{c}_{U}\right) ^{t}=\mathbf{\,}%
\left( \mathbf{c}_{U}^{\dagger }\mathbf{c}_{U}\right) ^{\ast }.
\end{equation}%
Hence we have that 
\begin{equation}
\left( \mathbf{cc}^{\dagger }\right) \mathbf{c-c}\left( \mathbf{cc}^{\dagger
}\right) ^{\ast }=\left( \mathbf{c}_{U}\mathbf{c}_{U}^{\dagger }\right) 
\mathbf{c}_{U}-\mathbf{c}_{U}\left( \mathbf{c}_{U}\mathbf{c}_{U}^{\dagger
}\right) ^{\ast }=0  \label{cu}
\end{equation}%
We can choose $\mathbb{U}$ to diagonalize the hermitian matrix $\mathbf{cc}%
^{\dagger }$ i.e.%
\begin{equation}
\mathbb{U}^{\dagger }\mathbf{cc}\mathbb{^{\dagger }U}=\left( 
\begin{array}{ccc}
n_{1} & 0 & 0 \\ 
0 & \ddots & 0 \\ 
0 & 0 & n_{r}%
\end{array}%
\right) =\mathbf{c}_{U}\mathbf{c}_{U}^{\dagger };\qquad n_{j}\geq 0\quad
1\leq j\leq r=2s.
\end{equation}%
$\lceil $The eigenvalues $\left\{ n_{j}\right\} $ are at least doubly
degenerate.$\rfloor $ Now using the relationships in eq. $\left( \ref{cu}%
\right) $ for this $\mathbf{c}_{U}$ we obtain%
\begin{equation}
\sum_{1\leq j\leq r}\left\{ n_{i}\delta _{ij}c_{Ujk}-c_{Uij}n_{j}\delta
_{jk}\right\} =0\Rightarrow \left( n_{i}-n_{k}\right) c_{Uik}=0;\qquad 1\leq
i,k\leq r
\end{equation}%
which shows that 
\begin{equation}
n_{k}\neq n_{i}\Rightarrow c_{Uik}=0;\qquad \qquad 1\leq i,k\leq r
\end{equation}%
Thus the matrix $\mathbf{c}_{U}$ breaks into submatrices $\left\{ \mathbb{C}%
_{\kappa };1\leq \kappa \leq p\right\} $ corresponding to groups of $%
n^{\prime }s$ that are equal. Consider the partitioning of $\left\{ 1,\ldots
,r\right\} $ into $p$ groups that correspond to equal $n^{\prime }s$ i.e. $%
\left\{ \left\{ \sigma \left( 1\right) ,\ldots ,\sigma \left( m_{1}\right)
\right\} ,\cdots ,\left\{ \sigma \left( m_{p-1}+1\right) ,\ldots \sigma
\left( r\right) \right\} \right\} $ where $\sigma $ is a permutation of $%
\left\{ 1,\ldots ,r\right\} $ and 
\begin{eqnarray}
n_{\sigma \left( 1\right) } &=&n_{\sigma \left( 2\right) }=\cdots =n_{\sigma
\left( m_{1}\right) }=\eta _{1}  \nonumber \\
&&\vdots  \nonumber \\
n_{\sigma \left( m_{p-1}+1\right) } &=&n_{\sigma \left( m_{p-1}+2\right)
}=\cdots =n_{\sigma \left( r\right) }=\eta _{p}
\end{eqnarray}%
then the matrix $\mathbf{c}_{U}$ has the block diagonal form 
\begin{equation}
\mathbf{c}_{U}=\left( 
\begin{array}{cccc}
\mathbb{C}_{1} & \mathbb{O} & \cdots & \mathbb{O} \\ 
\mathbb{O} & \ddots & \mathbb{O} & \vdots \\ 
\vdots & \mathbb{O} & \ddots & \mathbb{O} \\ 
\mathbb{O} & \cdots & \mathbb{O} & \mathbb{C}_{p}%
\end{array}%
\right)
\end{equation}%
and 
\begin{equation}
\mathbb{C}_{\kappa }\mathbb{C}_{\kappa }^{\dagger }=\eta _{\kappa }\mathbb{I}%
_{d\left( \kappa \right) };\quad 1\leq \kappa \leq p
\end{equation}%
where $d\left( \kappa \right) $ is the dimension of $\kappa ^{th}$-block.

If $\eta _{\kappa }=0$ then 
\begin{equation}
\sum_{k\in \left\{ \sigma \left( m_{\kappa -1}+1\right) ,\ldots \sigma
\left( m_{\kappa }\right) \right\} }c_{\kappa ik}c_{\kappa jk}^{\ast
}=0\qquad \forall i,j\in \left\{ \sigma \left( m_{\kappa -1}+1\right)
,\ldots ,\sigma \left( m_{\kappa }\right) \right\}
\end{equation}%
so in particular 
\begin{equation}
\sum_{k\in \left\{ \sigma \left( m_{\kappa -1}+1\right) ,\ldots \sigma
\left( m_{\kappa }\right) \right\} }\left| c_{\kappa jk}\right| ^{2}=0\qquad
\forall j\in \left\{ \sigma \left( m_{\kappa -1}+1\right) ,\ldots ,\sigma
\left( m_{\kappa }\right) \right\}
\end{equation}%
which implies that $c_{\kappa jk}=0\quad \forall j,k\in \left\{ \sigma
\left( m_{\kappa -1}+1\right) ,\ldots ,\sigma \left( m_{\kappa }\right)
\right\} $.

If $\eta _{\kappa }\neq 0$ then $\mathbb{B}_{\kappa }=\frac{1}{\sqrt{\eta
_{\kappa }}}\mathbb{C}_{\kappa }$ is a $d\left( \kappa \right) \times
d\left( \kappa \right) $ antisymmetric unitary matrix and 
\begin{eqnarray}
\mathbb{B}_{\kappa }^{\dagger }\mathbb{B}_{\kappa } &=&\mathbb{B}_{\kappa }%
\mathbb{B}_{\kappa }^{\dagger }=\mathbb{I}_{d\left( \kappa \right) } 
\nonumber \\
&\Rightarrow &\mathbb{B}_{\kappa }^{\ast }\mathbb{B}_{\kappa }=\mathbb{B}%
_{\kappa }\mathbb{B}_{\kappa }^{\ast }
\end{eqnarray}%
i.e. $\left[ \mathbb{B}_{\kappa },\mathbb{B}_{\kappa }^{\ast }\right] =0$
thus the matrices (\emph{not in general} self adjoint) 
\begin{eqnarray}
\mathbb{N}_{\kappa } &=&\frac{1}{2}\left\{ \mathbb{B}_{\kappa }+\mathbb{B}%
_{\kappa }^{\ast }\right\}  \nonumber \\
\mathbb{M}_{\kappa } &=&\frac{-i}{2}\left\{ \mathbb{B}_{\kappa }-\mathbb{B}%
_{\kappa }^{\ast }\right\}
\end{eqnarray}%
are real antisymmetric and \emph{commute }with each other so that 
\begin{equation}
\mathbb{B}_{\kappa }=\func{Re}\mathbb{B}_{\kappa }+i\func{Im}\mathbb{B}%
_{\kappa }.
\end{equation}
These matrices can be simultaneously transformed into the canonical forms 
\begin{equation}
\left( 
\begin{array}{cc}
\mathbf{0} & \mathbf{\alpha }_{\kappa } \\ 
-\mathbf{\alpha }_{\kappa } & \mathbf{0}%
\end{array}%
\right) \text{ and }\left( 
\begin{array}{cc}
\mathbf{0} & \mathbf{\beta }_{\kappa } \\ 
-\mathbf{\beta }_{\kappa } & \mathbf{0}%
\end{array}%
\right)
\end{equation}%
by a real orthogonal transformation $\mathbb{O}_{\kappa }$, where 
\begin{equation}
\mathbf{\alpha }_{\kappa }=\left( 
\begin{array}{ccc}
\alpha _{\kappa 1} & 0 & 0 \\ 
0 & \ddots & 0 \\ 
0 & 0 & \alpha _{\kappa d\left( \kappa \right) }%
\end{array}%
\right) \text{ and }\mathbf{\beta }_{\kappa }=\left( 
\begin{array}{ccc}
\beta _{\kappa 1} & 0 & 0 \\ 
0 & \ddots & 0 \\ 
0 & 0 & \beta _{\kappa d\left( \kappa \right) }%
\end{array}%
\right)
\end{equation}%
The same transformation thus acts on $\mathbb{B}_{\kappa }$ to produce 
\begin{equation}
\mathbb{O}_{\kappa }^{t}\mathbb{B}_{\kappa }\mathbb{O}_{\kappa }=\mathbb{O}%
_{\kappa }^{t}\left\{ \mathbb{N}_{\kappa }+i\mathbb{M}_{\kappa }\right\} 
\mathbb{O}_{\kappa }=\left( 
\begin{array}{cc}
\mathbf{0} & \mathbf{\alpha }_{\kappa }+i\mathbf{\beta }_{\kappa } \\ 
-\left( \mathbf{\alpha }_{\kappa }+i\mathbf{\beta }_{\kappa }\right) & 
\mathbf{0}%
\end{array}%
\right) =\left( 
\begin{array}{cc}
\mathbf{0} & e^{i\mathbf{\theta }_{\kappa }} \\ 
-e^{i\mathbf{\theta }_{\kappa }} & \mathbf{0}%
\end{array}%
\right)
\end{equation}%
and thus%
\begin{equation}
\mathbb{O}_{\kappa }^{t}\mathbb{C}_{\kappa }\mathbb{O}_{\kappa }=\eta
_{\kappa }^{\frac{1}{2}}\left( 
\begin{array}{cc}
\mathbf{0} & e^{i\mathbf{\theta }_{\kappa }} \\ 
-e^{i\mathbf{\theta }_{\kappa }} & \mathbf{0}%
\end{array}%
\right) =\left( 
\begin{array}{cc}
\mathbf{0} & \eta _{\kappa }^{\frac{1}{2}}e^{i\mathbf{\theta }_{\kappa }} \\ 
-\eta _{\kappa }^{\frac{1}{2}}e^{i\mathbf{\theta }_{\kappa }} & \mathbf{0}%
\end{array}%
\right) =\left( 
\begin{array}{cc}
\mathbf{0} & \mathbf{\zeta }_{\kappa } \\ 
-\mathbf{\zeta }_{\kappa } & \mathbf{0}%
\end{array}%
\right)
\end{equation}%
(which is valid even when $\eta _{\kappa }=0$) $\lceil $Note that $\mathbf{%
\theta }_{\kappa }\equiv \left( \theta _{\kappa 1},\ldots ,\theta _{\kappa
d\left( \kappa \right) }\right) $ and $\theta _{\kappa i}\neq \theta
_{\kappa j}$ in general$\rfloor $ where 
\begin{equation}
\mathbf{\zeta }_{\kappa }=\left( 
\begin{array}{ccc}
\zeta _{\kappa 1} & 0 & 0 \\ 
0 & \ddots & 0 \\ 
0 & 0 & \zeta _{\kappa d\left( \kappa \right) }%
\end{array}%
\right) .
\end{equation}%
This allows us to define 
\begin{equation}
\mathbb{O=}\left( 
\begin{array}{ccc}
\mathbb{O}_{1} & \mathbf{0} & \mathbf{0} \\ 
\mathbf{0} & \ddots & \mathbf{0} \\ 
\mathbf{0} & \mathbf{0} & \mathbb{O}_{p}%
\end{array}%
\right)
\end{equation}%
and see that 
\begin{equation}
\mathbb{O}^{t}\mathbf{c}_{U}\mathbb{O}=\left( 
\begin{array}{ccc}
\begin{array}{cc}
\mathbf{0} & \mathbf{\zeta }_{1} \\ 
-\mathbf{\zeta }_{1} & \mathbf{0}%
\end{array}
& \mathbf{0} & \mathbf{0} \\ 
\mathbf{0} & \ddots & \mathbf{0} \\ 
\mathbf{0} & \mathbf{0} & 
\begin{array}{cc}
\mathbf{0} & \mathbf{\zeta }_{p} \\ 
-\mathbf{\zeta }_{p} & \mathbf{0}%
\end{array}%
\end{array}%
\right)  \label{cor1}
\end{equation}%
and by rearranging the columns of $\mathbb{O}$ to produce $\widetilde{%
\mathbb{O}}$ we obtain 
\begin{equation}
\widetilde{\mathbb{O}}^{t}\mathbf{c}_{U}\widetilde{\mathbb{O}}=\left( 
\begin{array}{cc}
\mathbf{0} & \mathbf{\zeta } \\ 
-\mathbf{\zeta } & \mathbf{0}%
\end{array}%
\right)
\end{equation}%
and hence using $\mathbf{c}_{U}=\mathbb{U}^{\dagger }\mathbf{c}\mathbb{U}%
^{\ast }$ we conclude that 
\begin{equation}
\mathbb{T}\mathbf{c}\mathbb{T}^{t}=\widetilde{\mathbb{O}}^{t}\mathbb{U}%
^{\dagger }\mathbf{c}\mathbb{U}^{\ast }\widetilde{\mathbb{O}}=\left( 
\begin{array}{cc}
\mathbf{0} & \mathbf{\zeta } \\ 
-\mathbf{\zeta } & \mathbf{0}%
\end{array}%
\right)
\end{equation}%
and define $\mathbb{T}$ as the unitary matrix 
\begin{equation}
\mathbb{T}=\widetilde{\mathbb{O}}^{t}\mathbb{U}^{\dagger }
\end{equation}%
where $\mathbb{TT}^{\dagger }=\mathbb{T^{\dagger }T=I}_{r}$.

\begin{corollary}
If $\mathbf{c}$ is a $2s\times 2s$ complex antisymmetric matrix then there
exists a unitary matrix $\mathbb{W}$ such that 
\begin{equation}
\mathbb{W}\mathbf{c}\mathbb{W}^{t}\mathbb{=}\left( 
\begin{array}{cc}
\mathbf{0} & \mathbf{n}^{\frac{1}{2}} \\ 
-\mathbf{n}^{\frac{1}{2}} & \mathbf{0}%
\end{array}%
\right)
\end{equation}%
where 
\begin{equation}
\mathbf{n}^{\frac{1}{2}}=\left( 
\begin{array}{ccc}
n_{1}^{\frac{1}{2}} & 0 & 0 \\ 
0 & \ddots & 0 \\ 
0 & 0 & n_{s}^{\frac{1}{2}}%
\end{array}%
\right) ;\qquad n_{j}^{\frac{1}{2}}\in \mathbb{R},\quad 1\leq j\leq s
\end{equation}%
and $\left\{ n_{j}\right\} $ are the eigenvalues of $\mathbf{cc}^{\dagger }$.
\end{corollary}
\end{proof}
\end{theorem}

\begin{proof}
Using eq. $\left( \ref{cor1}\right) $ we can define $\mathbb{W}$ as 
\begin{equation}
\mathbb{W=}\left( 
\begin{array}{cc}
\mathbf{0} & e^{-i\frac{\mathbf{\theta }}{2}} \\ 
e^{-i\frac{\mathbf{\theta }}{2}} & \mathbf{0}%
\end{array}%
\right) \mathbb{T}
\end{equation}%
where 
\begin{equation}
\mathbb{W}^{\dagger }\mathbb{W=WW}^{\dagger }\mathbb{=I}_{2s}
\end{equation}%
and 
\begin{equation}
\mathbf{c}=\mathbb{W}^{\dagger }\left( 
\begin{array}{cc}
\mathbf{0} & \mathbf{n}^{\frac{1}{2}} \\ 
-\mathbf{n}^{\frac{1}{2}} & \mathbf{0}%
\end{array}%
\right) \mathbb{W}^{\ast }=\mathbb{T}^{\dagger }\left( 
\begin{array}{cc}
\mathbf{0} & \mathbf{\zeta } \\ 
-\mathbf{\zeta } & \mathbf{0}%
\end{array}%
\right) \mathbb{T}^{\ast }  \label{cex}
\end{equation}
\end{proof}

Eq $\left( \ref{cex}\right) $ shows that the canonical NGSOs $\left\{
\left\vert \psi _{k}\right\rangle \right\} $ are represented by the columns
of the matrix $\mathbb{W}^{\dagger }$ in the basis $\left\{ \left\vert
\varphi _{k}\right\rangle \right\} $ giving 
\begin{eqnarray}
W\left( \left\vert \mathbf{\psi }\right\rangle \right) &=&\left( \left\vert 
\mathbf{\varphi }\right\rangle \right)  \nonumber \\
W^{\dagger }\left( \left\vert \mathbf{\varphi }\right\rangle \right)
&=&\left( \left\vert \mathbf{\varphi }\right\rangle \right) \mathbb{W}%
^{\dagger }=\left( \left\vert \mathbf{\psi }\right\rangle \right)  \nonumber
\\
W^{\dagger }\left( \left\vert \mathbf{\varphi }\right\rangle \right)
&=&\left( \left\vert \mathbf{\varphi }\right\rangle \right) \left( 
\begin{array}{ccc}
\mathbf{w}_{1}^{\dagger } & \cdots & \mathbf{w}_{2s}^{\dagger }%
\end{array}%
\right)
\end{eqnarray}%
where 
\begin{eqnarray}
\left\vert g\right\rangle &=&\frac{1}{2}\sum_{1\leq k_{1},k_{2}\leq
r}c_{k_{1}k_{2}}\left\vert \varphi _{k_{1}}\varphi _{k_{2}}\right\rangle =%
\frac{1}{2}\sum_{1\leq k_{1},k_{2}\leq r}\left( \mathbb{W}^{\dagger }\left( 
\begin{array}{cc}
\mathbf{0} & \mathbf{n}^{\frac{1}{2}} \\ 
-\mathbf{n}^{\frac{1}{2}} & \mathbf{0}%
\end{array}%
\right) \mathbb{W}^{\ast }\right) _{k_{1}k_{2}}\left\vert \varphi
_{k_{1}}\varphi _{k_{2}}\right\rangle  \nonumber \\
&=&\sum_{1\leq k_{1}\leq s}n_{k_{1}}^{\frac{1}{2}}\left\vert W^{\dagger
}\varphi _{k_{1}}W^{\dagger }\varphi _{k_{1}+s}\right\rangle =\sum_{1\leq
k_{1}\leq s}n_{k_{1}}^{\frac{1}{2}}\left\vert \psi _{k_{1}}\psi
_{k_{1}+s}\right\rangle  \nonumber \\
&=&\sum_{1\leq k_{1}\leq s}n_{k_{1}}^{\frac{1}{2}}\left\vert e^{i\frac{%
\mathbf{\theta }_{k_{1}}}{2}}T^{\dagger }\varphi _{k_{1}}e^{i\frac{\mathbf{%
\theta }_{k_{1}}}{2}}T^{\dagger }\varphi _{k_{1}+s}\right\rangle
=\sum_{1\leq k_{1}\leq s}n_{k_{1}}^{\frac{1}{2}}e^{i\theta }\left\vert
T^{\dagger }\varphi _{k_{1}}T^{\dagger }\varphi _{k_{1}+s}\right\rangle
=\sum_{1\leq k_{1}\leq s}n_{k_{1}}^{\frac{1}{2}}e^{i\theta }\left\vert \chi
_{k_{1}}\chi _{k_{1}+s}\right\rangle
\end{eqnarray}%
and 
\begin{eqnarray}
T\left( \left\vert \mathbf{\chi }\right\rangle \right) &=&\left( \left\vert 
\mathbf{\varphi }\right\rangle \right)  \nonumber \\
T^{\dagger }\left( \left\vert \mathbf{\chi }\right\rangle \right) &=&\left(
\left\vert \mathbf{\varphi }\right\rangle \right) \mathbb{T}^{\dagger
}=\left( \left\vert \mathbf{\chi }\right\rangle \right)  \nonumber \\
T^{\dagger }\left( \left\vert \mathbf{\varphi }\right\rangle \right)
&=&\left( \left\vert \mathbf{\varphi }\right\rangle \right) \left( 
\begin{array}{ccc}
\mathbf{t}_{1}^{\dagger } & \cdots & \mathbf{t}_{2s}^{\dagger }%
\end{array}%
\right)
\end{eqnarray}

\section{Construction of $\mathbb{O}_{\protect\kappa }$ and Phase Factors}

We first form the self adjoint real and imaginary parts of $\mathbb{B}$ 
\begin{eqnarray}
\mathbb{B}_{R} &=&\func{Re}\mathbb{B}=\frac{1}{2}\left\{ \mathbb{B+B}%
^{\dagger }\right\}   \nonumber \\
\mathbb{B}_{I} &=&\func{Im}\mathbb{B}=-\frac{i}{2}\left\{ \mathbb{B-B}%
^{\dagger }\right\} 
\end{eqnarray}%
where we drop the subscript $\kappa $ which plays no role in the
considerations of this section. We know that 
\begin{equation}
\mathbb{B}_{R}=-\mathbb{B}_{R}^{t}\text{ and }\mathbb{B}_{I}=-\mathbb{B}%
_{I}^{t}
\end{equation}%
and 
\begin{eqnarray}
\mathbb{BB}^{\dagger } &=&\mathbb{B}^{\dagger }\mathbb{B=I}  \nonumber \\
\mathbb{BB}^{\ast } &=&\mathbb{B^{\ast }B}=-\mathbb{I}  \nonumber \\
\mathbb{B}^{-1} &=&-\mathbb{B^{\ast }}\text{ }
\end{eqnarray}%
so that 
\begin{equation}
\left[ \mathbb{B},\mathbb{B}^{\dagger }\right] =\mathbb{O}
\end{equation}%
and $\mathbb{B}$ is a \emph{normal }operator, thus its self adjoint real and
imaginary parts $\left\{ \mathbb{B}_{R},\mathbb{B}_{I}\right\} $ commute
(that are in general complex) and can be simultaneously diagonalized to
bring $\mathbb{B}$  to diagonal form. One can show that 
\begin{eqnarray}
\mathbb{B}_{R} &=&i\mathbb{M}  \nonumber \\
\mathbb{B}_{I} &=&-i\mathbb{N}
\end{eqnarray}%
where $\mathbb{N}$ and $\mathbb{M}$ are the real and imaginary parts of $%
\mathbb{B}$ (that are real) i.e. 
\begin{eqnarray}
\mathbb{N} &=&\frac{1}{2}\left\{ \mathbb{B+B}^{\ast }\right\}   \nonumber \\
\mathbb{M} &=&-\frac{i}{2}\left\{ \mathbb{B-B}^{\ast }\right\} 
\end{eqnarray}%
which also commute with each other 
\begin{equation}
\left[ \mathbb{M},\mathbb{N}\right] =\mathbb{O}
\end{equation}%
so that 
\begin{equation}
\mathbb{B=B}_{R}+i\mathbb{B}_{I}=\mathbb{N}+i\mathbb{M}
\end{equation}%
We can simultaneously diagonalize the (in general complex) antisymmetric
self adjoint matrices $\mathbb{B}_{R}$ and $\mathbb{B}_{I}$ and hence $%
\mathbb{B}$ with a unitary matrix $\mathbb{V}$ to obtain 
\begin{equation}
\mathbb{B}=\left( \mathbf{v}_{1}\mathbf{v}_{1}^{\ast }\mathbf{\cdots v}%
_{s^{\prime }}\mathbf{v}_{s^{\prime }}^{\ast }\right) \left( 
\begin{array}{ccccc}
-ie^{i\theta _{1}} & 0 & 0 & 0 & 0 \\ 
0 & ie^{i\theta _{1}} & 0 & 0 & 0 \\ 
0 & 0 & \ddots  & 0 & 0 \\ 
0 & 0 & 0 & -ie^{i\theta _{s^{\prime }}} & 0 \\ 
0 & 0 & 0 & 0 & ie^{i\theta _{s^{\prime }}}%
\end{array}%
\right) \left( 
\begin{array}{c}
\mathbf{v}_{1}^{\ast } \\ 
\mathbf{v}_{1} \\ 
\mathbf{\vdots } \\ 
\mathbf{v}_{s^{\prime }}^{\ast } \\ 
\mathbf{v}_{s^{\prime }}%
\end{array}%
\right) 
\end{equation}%
$\lceil $In eigen vector form%
\begin{eqnarray}
\mathbb{B}\mathbf{v}_{1}^{\ast } &=&-ie^{i\theta _{1}}\mathbf{v}_{1}^{\ast }
\\
\mathbb{B}\mathbf{v}_{1} &=&ie^{i\theta _{1}}\mathbf{v}_{1}
\end{eqnarray}%
\begin{eqnarray}
\mathbb{B} &=&\left( 
\begin{array}{cc}
\mathbf{v} & \overline{\mathbf{v}}%
\end{array}%
\right) \left( 
\begin{array}{cc}
ie^{i\mathbf{\theta }} & 0 \\ 
0 & -ie^{i\mathbf{\theta }}%
\end{array}%
\right) \left( 
\begin{array}{c}
\overline{\mathbf{v}} \\ 
\mathbf{v}%
\end{array}%
\right)  \\
&=&i\left( 
\begin{array}{cc}
e^{i\frac{\mathbf{\theta }}{2}}\mathbf{v} & e^{i\frac{\mathbf{\theta }}{2}}%
\overline{\mathbf{v}}%
\end{array}%
\right) \left( 
\begin{array}{cc}
I & 0 \\ 
0 & -I%
\end{array}%
\right) \left( 
\begin{array}{c}
e^{i\frac{\mathbf{\theta }}{2}}\overline{\mathbf{v}} \\ 
e^{i\frac{\mathbf{\theta }}{2}}\mathbf{v}%
\end{array}%
\right) ,
\end{eqnarray}

which is not a unitary transformation of the basis$\rfloor $

and hence the \emph{real} orthogonal matrix $\mathbb{O}$ by defining 
\begin{eqnarray}
\mathbf{o}_{j} &=&\frac{1}{\sqrt{2}}\left( \mathbf{v}_{j}+\mathbf{v}%
_{j}^{\ast }\right) =\sqrt{2}\func{Re}\mathbf{v}_{j}  \nonumber \\
\mathbf{o}_{j+s^{\prime }} &=&\frac{-i}{\sqrt{2}}\left( \mathbf{v}_{j}-%
\mathbf{v}_{j}^{\ast }\right) =\sqrt{2}\func{Im}\mathbf{v}_{j}
\end{eqnarray}%
so that 
\begin{eqnarray}
\mathbb{B} &=&i\left( \mathbf{o}_{1}\mathbf{o}_{1+s^{\prime }}\mathbf{\cdots
o}_{s^{\prime }}\mathbf{o}_{s^{\prime }+s^{\prime }}\right) \left( 
\begin{array}{ccccc}
0 & e^{i\theta _{1}} & 0 & 0 & 0 \\ 
-e^{i\theta _{1}} & 0 & 0 & 0 & 0 \\ 
0 & 0 & \ddots  & 0 & 0 \\ 
0 & 0 & 0 & 0 & e^{i\theta _{s^{\prime }}} \\ 
0 & 0 & 0 & -e^{i\theta _{s^{\prime }}} & 0%
\end{array}%
\right) \left( 
\begin{array}{c}
\mathbf{o}_{1} \\ 
\mathbf{o}_{1+s^{\prime }} \\ 
\mathbf{\vdots } \\ 
\mathbf{o}_{s^{\prime }} \\ 
\mathbf{o}_{s^{\prime }+s^{\prime }}%
\end{array}%
\right)   \nonumber \\
\mathbb{B} &=&i\left( \mathbf{\tilde{o}}_{1}\mathbf{\tilde{o}}_{1+s^{\prime
}}\mathbf{\cdots \tilde{o}}_{\mathbf{s^{\prime }}}\mathbf{\tilde{o}}%
_{s^{\prime }+s^{\prime }}\right) \left( 
\begin{array}{ccccc}
0 & 1 & 0 & 0 & 0 \\ 
-1 & 0 & 0 & 0 & 0 \\ 
0 & 0 & \ddots  & 0 & 0 \\ 
0 & 0 & 0 & 0 & 1 \\ 
0 & 0 & 0 & -1 & 0%
\end{array}%
\right) \left( 
\begin{array}{c}
\mathbf{\tilde{o}} \\ 
\mathbf{\tilde{o}}_{1+s^{\prime }} \\ 
\mathbf{\vdots } \\ 
\mathbf{\tilde{o}}_{s^{\prime }} \\ 
\mathbf{\tilde{o}}_{s^{\prime }+s^{\prime }}%
\end{array}%
\right)   \label{reex}
\end{eqnarray}%
where 
\begin{equation}
\left( \mathbf{\tilde{o}}_{1}\mathbf{\tilde{o}}_{1+s^{\prime }}\mathbf{%
\cdots \tilde{o}}_{\mathbf{s^{\prime }}}\mathbf{\tilde{o}}_{s^{\prime
}+s^{\prime }}\right) =\left( \mathbf{o}_{1}\mathbf{o}_{1+s^{\prime }}%
\mathbf{\cdots o}_{s^{\prime }}\mathbf{o}_{s^{\prime }+s^{\prime }}\right)
\left( 
\begin{array}{ccccc}
0 & e^{i\frac{\theta _{1}}{2}} & 0 & 0 & 0 \\ 
-e^{i\frac{\theta _{1}}{2}} & 0 & 0 & 0 & 0 \\ 
0 & 0 & \ddots  & 0 & 0 \\ 
0 & 0 & 0 & 0 & e^{i\frac{\theta _{s^{\prime }}}{2}} \\ 
0 & 0 & 0 & -e^{i\frac{\theta _{s^{\prime }}}{2}} & 0%
\end{array}%
\right) .
\end{equation}%
This last relationship gives 
\begin{equation}
\mathbb{B}=i\mathbf{\hat{o}}\left( 
\begin{array}{ccccc}
0 & 1 & 0 & 0 & 0 \\ 
-1 & 0 & 0 & 0 & 0 \\ 
0 & 0 & \ddots  & 0 & 0 \\ 
0 & 0 & 0 & 0 & 1 \\ 
0 & 0 & 0 & -1 & 0%
\end{array}%
\right) \mathbf{\hat{o}}^{t}=i\mathbf{\hat{o}}\left( 
\begin{array}{cc}
\mathbb{O} & \mathbb{I} \\ 
-\mathbb{I} & \mathbb{O}%
\end{array}%
\right) \mathbf{\hat{o}}^{t}.
\end{equation}%
Thus 
\begin{equation}
\mathbb{B}=i\sum_{1\leq j\leq s^{\prime }}e^{i\theta _{j}}\left( \left\vert 
\mathbf{o}_{j}\right\rangle \left\langle \mathbf{o}_{j+s^{\prime
}}\right\vert -\left\vert \mathbf{o}_{j+s^{\prime }}\right\rangle
\left\langle \mathbf{o}_{j}\right\vert \right) 
\end{equation}%
and 
\begin{eqnarray}
\mathbb{B} &=&i\sum_{1\leq j\leq s^{\prime }}\left( \left\vert e^{i\frac{%
\theta _{j}}{2}}\mathbf{o}_{j}\right\rangle \left\langle e^{-i\frac{\theta
_{j}}{2}}\mathbf{o}_{j+s^{\prime }}\right\vert -\left\vert e^{i\frac{\theta
_{j}}{2}}\mathbf{o}_{j+s^{\prime }}\right\rangle \left\langle e^{-i\frac{%
\theta _{j}}{2}}\mathbf{o}_{j}\right\vert \right)  \\
&=&i\sum_{1\leq j\leq s^{\prime }}\left( \left\vert \mathbf{\tilde{o}}%
_{j}\right\rangle \left\langle \mathbf{\tilde{o}}_{j+s^{\prime }}\right\vert
-\left\vert \mathbf{\tilde{o}}_{j+s^{\prime }}\right\rangle \left\langle 
\mathbf{\tilde{o}}_{j}\right\vert \right) .
\end{eqnarray}%
It should be noted that 
\begin{equation}
\mathbb{B}=\left( 
\begin{array}{cc}
\mathbf{w} & \overline{\mathbf{w}}%
\end{array}%
\right) \left( 
\begin{array}{cc}
\mathbb{I} & 0 \\ 
0 & -\mathbb{I}%
\end{array}%
\right) \left( 
\begin{array}{c}
\overline{\mathbf{w}} \\ 
\mathbf{w}%
\end{array}%
\right) ,
\end{equation}%
where 
\begin{equation}
\mathbf{w}=e^{-i\left( \mathbf{\theta }+\frac{\pi }{2}\right) }\mathbf{v.}
\end{equation}

\end{document}
